% Generated draft: requires mathematical and source review.
\documentclass[11pt]{article}
\usepackage[margin=25mm]{geometry}
\usepackage{fontspec}
\setmainfont{Latin Modern Roman}
\usepackage{amsmath,amssymb,mathtools,unicode-math}
\setmathfont{Latin Modern Math}
\usepackage{enumitem,needspace,xurl,hyperref}
\setlength{\parindent}{0pt}
\setlength{\parskip}{5pt}
\setlength{\emergencystretch}{3em}
\newcommand{\N}{\mathbb{N}}
\newcommand{\Z}{\mathbb{Z}}
\newcommand{\Q}{\mathbb{Q}}
\newcommand{\R}{\mathbb{R}}
\newcommand{\C}{\mathbb{C}}
\newcommand{\F}{\mathbb{F}}
\newcommand{\A}{\mathbb{A}}
\newcommand{\PP}{\mathbb{P}}
\newcommand{\E}{\mathbb{E}}
\newcommand{\eps}{\varepsilon}
\newcommand{\dd}{\,\mathrm{d}}
\begin{document}
\section*{UnsolvedMath ID: 701751}
\textbf{Source problem code:} catalogue entry
\subsection{Definitions and mathematical statement}
For a subgroup \(H\leq G\), let
\[
N_G(H)=\{g\in G:g^{-1}Hg=H\}.
\]
The subgroup \(H\) is normal in \(N_G(H)\), so the quotient \(N_G(H)/H\) is defined. A group \(G\) is a \emph{Shunkov group} if, for every finite subgroup \(H\leq G\), every prime \(p\), and all \(x,u\in N_G(H)/H\) with \(x\) of order \(p\),
\[
\bigl|\langle x,u^{-1}xu\rangle\bigr|<\infty.
\]
Thus any two conjugate elements of prime order in each such quotient generate a finite subgroup. A group is \emph{periodic} if every element has finite order. An \emph{involution} is an element of order \(2\).

The supplied problem asks whether the following results from finite group theory remain true in the class of ``(periodic) Shunkov groups'':
\begin{enumerate}
\item[(a)] The Baer--Suzuki theorem, with the source referring to Archive problem 11.11.
\item[(b)] The Burnside--Brauer--Suzuki theorem, described in the source as asserting the existence of a normal section of order \(2\) in a group with a nontrivial Sylow \(2\)-subgroup containing exactly one involution; the source refers to problem 4.75.
\item[(c)] Glauberman's \(Z^*\)-theorem, with the source referring to problem 10.62 and Archive problem 11.13.
\end{enumerate}

\emph{Unresolved qualifications.} The parentheses around ``periodic'' do not specify whether periodicity is required throughout, whether separate periodic and unrestricted versions are intended, or whether it is an optional strengthening. The cited theorem statements and referenced notebook entries are not included in the supplied record. Consequently, their precise hypotheses, conclusions, and intended interpretations for infinite groups cannot be recovered from this record alone. In particular, part (b) does not specify the intended infinite-group meaning of a Sylow \(2\)-subgroup or which subgroups defining the normal section must be normal in the ambient group. The named extension questions above therefore retain these unresolved qualifications rather than supplying additional hypotheses or theorem formulations.
\subsection{Short English statement}
A Shunkov group satisfies a finiteness condition on pairs of conjugate prime-order elements in normalizer quotients of finite subgroups. The problem asks whether three named finite-group theorems extend to such groups, possibly under periodicity. The precise scope of periodicity and the full intended theorem formulations require clarification from the cited notebook entries.
\subsection{Sources}
\begin{itemize}
\item Local UnsolvedMath registry, supplied complete source record: permanent ID \(701751\), exact source problem code \texttt{catalogue entry}, title \emph{Kourovka Notebook Problem 20.93}. The registry attributes the problem to A. I. Sozutov, proposed in 2022, and labels it open; this label is supplied metadata, not an independently verified mathematical assessment.
\item Supplied bibliographic attribution: \emph{The Kourovka Notebook (Unsolved Problems in Group Theory)}, 20th issue (2022), original problem 20.93; also identified with the 21st issue edition (2026), \texttt{arXiv:1401.0300v48}. The collection metadata names E. I. Khukhro and V. D. Mazurov as editors.
\item Supplied internal references: Archive problem 11.11 for part (a); problem 4.75 for part (b); problem 10.62 and Archive problem 11.13 for part (c). Their contents are not supplied.
\item The supplied registry record was created and updated on \texttt{2026-10-16T14:00:00Z}, with upstream revision \texttt{d371ee917ac9647fbb06cd6fc529deec6f3dda82}. These are record provenance timestamps, not a supplied literature-verification date. No source was independently consulted for this draft.
\end{itemize}
\end{document}
